\documentclass[11pt,a4paper]{article}
% Wspólny szablon dla kart supermocy (jedna karta = jedna strona A4)
% Użycie w karcie: \documentclass[11pt,a4paper]{article} % Wspólny szablon dla kart supermocy (jedna karta = jedna strona A4)
% Użycie w karcie: \documentclass[11pt,a4paper]{article} % Wspólny szablon dla kart supermocy (jedna karta = jedna strona A4)
% Użycie w karcie: \documentclass[11pt,a4paper]{article} \input{../preambula} \begin{document}

% Silnik tectonic (XeTeX): polskie znaki przez fontspec, nie przez inputenc/T1
\usepackage{fontspec}
\setmainfont{Arial}
\usepackage[polish]{babel}
\usepackage[a4paper,margin=1.5cm]{geometry}
\usepackage{amsmath,amssymb}
\usepackage{tikz}
\usetikzlibrary{arrows.meta,positioning,calc,shapes.geometric}
\usepackage[most]{tcolorbox}
\usepackage{enumitem}
\usepackage{booktabs}
\usepackage{xcolor}

\pagestyle{empty}
\setlength{\parindent}{0pt}
\setlength{\parskip}{0.4em}

\definecolor{cel}{HTML}{1F4E79}
\definecolor{regula}{HTML}{E8F1FA}
\definecolor{pulapka}{HTML}{FDECEA}
\definecolor{okres}{HTML}{C0392B}
\definecolor{skonczony}{HTML}{1E8449}

% Ramka na regułę
\newtcolorbox{regula}[1][]{
  colback=regula, colframe=cel, boxrule=0.8pt, arc=2pt,
  left=6pt, right=6pt, top=4pt, bottom=4pt, title=#1,
  fonttitle=\bfseries, coltitle=white
}

% Ramka na pułapkę z kartkówki
\newtcolorbox{pulapka}[1][]{
  colback=pulapka, colframe=okres, boxrule=0.8pt, arc=2pt,
  left=6pt, right=6pt, top=4pt, bottom=4pt, title=#1,
  fonttitle=\bfseries, coltitle=white
}

% Nagłówek karty: numer, tytuł, cel
\newcommand{\naglowekKarty}[3]{%
  {\LARGE\bfseries\color{cel} Karta #1: #2}\\[2pt]
  {\large Cel: #3}\par
  \vspace{0.3em}\noindent\rule{\linewidth}{1pt}\par
}

% Pole na obliczenia ucznia: ramka z delikatną kratką, argument = wysokość
\newcommand{\miejsce}[1]{%
  \par\vspace{0.2em}\noindent
  \begin{tikzpicture}
    \draw[step=0.5cm, gray!25, very thin] (0,0) grid (\textwidth,#1);
    \draw[gray!60, rounded corners=2pt] (0,0) rectangle (\textwidth,#1);
    \node[anchor=north west, font=\scriptsize, text=gray] at (0.05,#1) {obliczenia};
  \end{tikzpicture}\par}
 \begin{document}

% Silnik tectonic (XeTeX): polskie znaki przez fontspec, nie przez inputenc/T1
\usepackage{fontspec}
\setmainfont{Arial}
\usepackage[polish]{babel}
\usepackage[a4paper,margin=1.5cm]{geometry}
\usepackage{amsmath,amssymb}
\usepackage{tikz}
\usetikzlibrary{arrows.meta,positioning,calc,shapes.geometric}
\usepackage[most]{tcolorbox}
\usepackage{enumitem}
\usepackage{booktabs}
\usepackage{xcolor}

\pagestyle{empty}
\setlength{\parindent}{0pt}
\setlength{\parskip}{0.4em}

\definecolor{cel}{HTML}{1F4E79}
\definecolor{regula}{HTML}{E8F1FA}
\definecolor{pulapka}{HTML}{FDECEA}
\definecolor{okres}{HTML}{C0392B}
\definecolor{skonczony}{HTML}{1E8449}

% Ramka na regułę
\newtcolorbox{regula}[1][]{
  colback=regula, colframe=cel, boxrule=0.8pt, arc=2pt,
  left=6pt, right=6pt, top=4pt, bottom=4pt, title=#1,
  fonttitle=\bfseries, coltitle=white
}

% Ramka na pułapkę z kartkówki
\newtcolorbox{pulapka}[1][]{
  colback=pulapka, colframe=okres, boxrule=0.8pt, arc=2pt,
  left=6pt, right=6pt, top=4pt, bottom=4pt, title=#1,
  fonttitle=\bfseries, coltitle=white
}

% Nagłówek karty: numer, tytuł, cel
\newcommand{\naglowekKarty}[3]{%
  {\LARGE\bfseries\color{cel} Karta #1: #2}\\[2pt]
  {\large Cel: #3}\par
  \vspace{0.3em}\noindent\rule{\linewidth}{1pt}\par
}

% Pole na obliczenia ucznia: ramka z delikatną kratką, argument = wysokość
\newcommand{\miejsce}[1]{%
  \par\vspace{0.2em}\noindent
  \begin{tikzpicture}
    \draw[step=0.5cm, gray!25, very thin] (0,0) grid (\textwidth,#1);
    \draw[gray!60, rounded corners=2pt] (0,0) rectangle (\textwidth,#1);
    \node[anchor=north west, font=\scriptsize, text=gray] at (0.05,#1) {obliczenia};
  \end{tikzpicture}\par}
 \begin{document}

% Silnik tectonic (XeTeX): polskie znaki przez fontspec, nie przez inputenc/T1
\usepackage{fontspec}
\setmainfont{Arial}
\usepackage[polish]{babel}
\usepackage[a4paper,margin=1.5cm]{geometry}
\usepackage{amsmath,amssymb}
\usepackage{tikz}
\usetikzlibrary{arrows.meta,positioning,calc,shapes.geometric}
\usepackage[most]{tcolorbox}
\usepackage{enumitem}
\usepackage{booktabs}
\usepackage{xcolor}

\pagestyle{empty}
\setlength{\parindent}{0pt}
\setlength{\parskip}{0.4em}

\definecolor{cel}{HTML}{1F4E79}
\definecolor{regula}{HTML}{E8F1FA}
\definecolor{pulapka}{HTML}{FDECEA}
\definecolor{okres}{HTML}{C0392B}
\definecolor{skonczony}{HTML}{1E8449}

% Ramka na regułę
\newtcolorbox{regula}[1][]{
  colback=regula, colframe=cel, boxrule=0.8pt, arc=2pt,
  left=6pt, right=6pt, top=4pt, bottom=4pt, title=#1,
  fonttitle=\bfseries, coltitle=white
}

% Ramka na pułapkę z kartkówki
\newtcolorbox{pulapka}[1][]{
  colback=pulapka, colframe=okres, boxrule=0.8pt, arc=2pt,
  left=6pt, right=6pt, top=4pt, bottom=4pt, title=#1,
  fonttitle=\bfseries, coltitle=white
}

% Nagłówek karty: numer, tytuł, cel
\newcommand{\naglowekKarty}[3]{%
  {\LARGE\bfseries\color{cel} Karta #1: #2}\\[2pt]
  {\large Cel: #3}\par
  \vspace{0.3em}\noindent\rule{\linewidth}{1pt}\par
}

% Pole na obliczenia ucznia: ramka z delikatną kratką, argument = wysokość
\newcommand{\miejsce}[1]{%
  \par\vspace{0.2em}\noindent
  \begin{tikzpicture}
    \draw[step=0.5cm, gray!25, very thin] (0,0) grid (\textwidth,#1);
    \draw[gray!60, rounded corners=2pt] (0,0) rectangle (\textwidth,#1);
    \node[anchor=north west, font=\scriptsize, text=gray] at (0.05,#1) {obliczenia};
  \end{tikzpicture}\par}


\begin{document}

\naglowekKarty{2}{Skończone czy nieskończone?}{rozpoznać rozwinięcie dziesiętne ułamka bez dzielenia „na oko”}

\vspace{0.3em}
\begin{regula}[REGUŁA]
\textbf{Krok 1.} Skróć ułamek do postaci nieskracalnej.\\
\textbf{Krok 2.} Spójrz na mianownik $b$ po skróceniu:
$$\frac{a}{b}\ \text{(nieskracalny)}\qquad\Longrightarrow\qquad
\text{rozwinięcie \emph{skończone}} \iff b = 2^{m}\cdot 5^{n},\quad m,n\geq 0$$
Jeśli w mianowniku jest choćby jeden inny czynnik pierwszy (np. $3$, $7$, $11$), rozwinięcie jest \emph{nieskończone okresowe}.
\end{regula}

\vspace{0.1em}
\textbf{Dlaczego tak? Druga droga: dzielenie pod kreską.} Dzielenie kończy się, gdy reszta wynosi $0$. Reszt jest mniej niż dzielnik, więc jeśli dzielenie \emph{nie} kończy się, reszty się powtarzają, a z nimi cyfry wyniku: to jest \emph{okres}.

\begin{center}
\begin{tikzpicture}[scale=0.82, >=Stealth, node distance=2.55cm,
  box/.style={draw, rounded corners=2pt, align=center, minimum height=0.85cm, minimum width=2.3cm, inner sep=3pt}]
  % Wiersz A: 3 : 8 -> reszta 0 -> koniec
  \node[font=\bfseries, anchor=west] at (-4.6,0.85) {A) $\frac{3}{8}$};
  \node[box] (a1) at (-4.6,0) {$30:8=3$\\reszta $6$};
  \node[box, right=of a1] (a2) {$60:8=7$\\reszta $4$};
  \node[box, right=of a2] (a3) {$40:8=5$\\reszta $\mathbf{0}$};
  \node[box, fill=skonczony!20, draw=skonczony, right=of a3] (a4) {$0{,}375$\\\small{\textbf{KONIEC}}};
  \draw[->, thick] (a1) -- (a2);
  \draw[->, thick] (a2) -- (a3);
  \draw[->, thick] (a3) -- (a4);

  % Wiersz B: 2 : 3 -> reszta 2 się powtarza -> okres
  \node[font=\bfseries, anchor=west] at (-4.6,-1.95) {B) $\frac{2}{3}$};
  \node[box] (b1) at (-4.6,-2.8) {$20:3=6$\\reszta $\mathbf{2}$};
  \node[box, right=of b1] (b2) {$20:3=6$\\reszta $\mathbf{2}$};
  \node[box, fill=okres!12, draw=okres, right=of b2] (b3) {$0{,}(6)$\\\small{\textbf{ZNOWU } $2$ $\Rightarrow$ OKRES}};
  \draw[->, thick] (b1) -- (b2);
  \draw[->, thick] (b2) -- (b3);
  \draw[->, thick, color=okres] (b2.south) to[bend left=45] node[below, font=\small] {reszta $2$ wraca} (b1.south);
\end{tikzpicture}
\end{center}


\vspace{0.1em}
\textbf{Przykłady rozstrzygnięte regułą:}
\renewcommand{\arraystretch}{1.15}
\begin{center}
\begin{tabular}{@{}lllll@{}}
\toprule
Ułamek & Skracamy? & Mianownik po skróceniu & Rozwinięcie & Typ \\
\midrule
$\dfrac{3}{20}$ & nie & $20 = 2^{2}\cdot 5$ & $0{,}15$ & \textcolor{skonczony}{\textbf{skończone}} \\
$\dfrac{7}{40}$ & nie & $40 = 2^{3}\cdot 5$ & $0{,}175$ & \textcolor{skonczony}{\textbf{skończone}} \\
$\dfrac{6}{15}$ & tak, $:3$ & $\frac{2}{5}$, czyli $5$ & $0{,}4$ & \textcolor{skonczony}{\textbf{skończone}} \\
$\dfrac{4}{9}$ & nie & $9 = 3^{2}$ & $0{,}(4)$ & \textcolor{okres}{\textbf{nieskończone}} \\
$\dfrac{5}{12}$ & nie & $12 = 2^{2}\cdot 3$ & $0{,}41(6)$ & \textcolor{okres}{\textbf{nieskończone}} \\
\bottomrule
\end{tabular}
\end{center}

\begin{pulapka}[PUŁAPKI Z KARTKÓWKI]
\begin{enumerate}[leftmargin=*, itemsep=2pt]
  \item $\frac{6}{15}$ \emph{wygląda} na nieskończony, bo $15 = 3\cdot 5$. Po skróceniu przez $3$ zostaje $\frac{2}{5}$, więc jest skończony. \textbf{Najpierw skracaj!}
  \item $\frac{5}{12}$ ma w mianowniku dwójkę, ale trójka też wystarcza, żeby rozwinięcie było nieskończone. Jeden obcy czynnik psuje cały ułamek.
\end{enumerate}
\end{pulapka}

\begin{regula}[ZAMIANA NA DZIESIĘTNY (skończony)]
Mianownik zamień na $10^{k}$ przez dopisanie „przyjaciela”: $2\to 5$, $4\to 25$, $8\to 125$, $5\to 2$, $25\to 4$, $125\to 8$.\\
Np. $\displaystyle\frac{3}{8} = \frac{3\cdot 125}{8\cdot 125} = \frac{375}{1000} = 0{,}375$.
\end{regula}

\end{document}
